\expandafter\def\csname HistoricalNoteText16\endcsname{[Handschriftliche Bemerkung.] Am 31. Dez. schon fast 11 Grad. Vgl. S. 581 und 647. [S. 203 dieses Bandes.]}
\clearpage\SourceLeaf{204}%
\pdfbookmark[1]{Fortgesetzte Nachrichten. Ceres. Dezember 1801}{ceres180112}
\fancyhead[L]{\small [CERES.]}
\begin{center}
{\small FORTGESETZTE NACHRICHTEN ÜBER DEN LÄNGST VERMUTETEN NEUEN HAUPTPLANETEN UNSERES SONNENSYSTEMS.\par}
[CERES.]\par
\medskip\hrule\medskip
{\small Monatliche Correspondenz zur Beförderung der Erd- und Himmels-Kunde, herausgegeben von Freiherrn \textsc{von Zach}. Band III. Dezember 1801.\par}
\medskip\hrule
\end{center}\medskip
\noindent —\quad —\quad —\par
Die Untersuchungen und Berechnungen, die uns Dr.\ \textsc{Gauss} in Braunschweig vor Kurzem mitgeteilt hat, lassen uns sehr auf Hilfe und Erleichterung hoffen. Zugleich machen sie erneut sehr wahrscheinlich, dass das von \textsc{Piazzi} entdeckte neue Gestirn tatsächlich ein planetarischer Himmelskörper ist, der sich nach den \textsc{Kepler}'schen Gesetzen zwischen der Mars- und der Jupiterbahn bewegt.

Wir beeilen uns umso mehr, seine Berechnungen mitzuteilen, als sich seine neue elliptische Bahn beträchtlich von der elliptischen Bahn des Dr.\ \textsc{Burckhardt} und den beiden Kreisbahnen des Dr.\ \textsc{Olbers} und des Paters \textsc{Piazzi} unterscheidet. Diese Bahnen haben wir in unseren vorigen Heften zusammen mit den daraus vorausberechneten Orten des Planeten mitgeteilt. Ihre Abweichung von den \textsc{Gauss}'schen Positionen kann in diesem Monat 6 bis 7 Grad in geozentrischer Länge erreichen. Es ist daher wichtig, den beobachtenden Astronomen diese Bemerkungen so bald wie möglich mitzuteilen. Sie erfahren dadurch, dass sie das Gebiet am Himmel, in dem sie dieses neue und so schwer zu findende Gestirn suchen müssen, \textit{unbedingt} um weitere 6 bis 7 Grad nach Osten ausdehnen müssen. Denn 1) nach der Rechnung von Dr.\ \textsc{Gauss} liegen die Beobachtungen \textsc{Piazzi}s fast genau in der Mitte zwischen Perihel und Aphel. Die Doktoren \textsc{Burckhardt} und \textsc{Olbers} hatten dagegen angenommen, dass sie in der Nähe des einen oder des anderen liegen.

2) Die Exzentrizität der Bahn ist nach seiner Rechnung nicht ganz so gering, wie Pater \textsc{Piazzi} glaubt. Daher können sich die nach seinen elliptischen Elementen vorausberechneten Längen in diesem Monat leicht um 7 Grad\HistoricalNote{16} von den unter den Kreishypothesen berechneten unterscheiden.

3) Die \textsc{Gauss}'schen Elemente bestätigen die Beobachtungen \textsc{Piazzi}s vollkommen und räumen den Verdacht mangelnder Genauigkeit vollständig aus, den die Doktoren \textsc{Burckhardt} und \textsc{Olbers} bei der Prüfung anhand der Differenzen gehegt hatten.

All dies beweist die \textsc{Gauss}'sche Ellipse. Wie viel Vertrauen sie daher erwecken muss, werden Astronomen an der Genauigkeit erkennen, mit der sie sämtliche Beobachtungen \textsc{Piazzi}s darstellt. Dr.\ \textsc{Gauss} kam zu diesen Berechnungen durch einige Untersuchungen zur physischen Astronomie,
\SourceLeaf{205}%
die ihn zu einigen durchaus bedeutenden Ergänzungen der Theorie der Bestimmung der Bewegung von Himmelskörpern in Kegelschnitten jeder Art geführt haben. Er war so freundlich, uns einiges davon mitzuteilen. Wir hoffen, unsere astronomischen Leser ein andermal damit zu unterhalten, da uns eine Darstellung hier zu weit von unserem Gegenstand abbringen würde. Wir beschränken uns daher auf das, was sich unmittelbar auf die Untersuchung der Bahn der \textit{Ceres Ferdinandea} bezieht.

Zur Bestimmung der Bahn wählte Dr.\ \textsc{Gauss} zuerst die drei Beobachtungen vom 2. und 22. Januar und vom 11. Februar. Dabei übernahm er die Daten genau so, wie sie im Septemberheft der M.\ C. auf S.\ 280 angegeben sind. Mit einem eigenen Verfahren fand er schon beim vierten Versuch die folgenden Elemente:
\begin{center}
{\small\setlength{\tabcolsep}{2pt}
\begin{tabular}{l r}
Aphel $\ .\ .\ .\ .\ .\ .\ .\ .\ .\ .$&$330^0\ 14'\ 33''$\\
$\AscendingNode\ .\ .\ .\ .\ .\ .\ .\ .\ .\ .\ .\ .\ .$&$81\quad8\ 50$\\
Bahnneigung $\ .\ .\ .\ .\ .\ .\ .\ .$&$10\ 32\ 19$\\
Logarithmus der großen Halbachse $\ .\ .$&$0.4381058$\\
Exzentrizität $\ .\ .\ .\ .\ .\ .\ .\ .\ .\ .$&$0.0832836$\\
\multicolumn{2}{l}{Epoche: 31. Dez. 1800 am mittleren Mittag in Palermo}\\
mittlere heliozentrische Länge: $\ .\ .\ .$&$77^0\ 54'\ 29''$\\[4pt]
\multicolumn{2}{c}{Hieraus folgt:}\\[3pt]
größte Mittelpunktsgleichung $\ .\ .\ .$&$9^0\ 32'\ 57''$\\
große Halbachse $\ .\ .\ .\ .\ .\ .\ .\ .$&$2.74226$\\
siderische Umlaufzeit $\ .\ .\ .\ .\ .\ .$&$1658\frac23$ Tage\\
mittlere tägliche siderische Bewegung&$781''.355$
\end{tabular}}
\end{center}
Diese Elemente stellen die beiden äußeren Beobachtungen genau dar, die mittlere mit einem Fehler von $2''$ in Länge und Breite. Aberration und Präzession wurden angemessen berücksichtigt.

Nach diesem ersten gelungenen Versuch nahm Dr.\ \textsc{Gauss} eine zweite Berechnung dieser Bahn vor. Er ließ die Beobachtungen vom 2. und 22. Januar weg und wählte stattdessen die vom 1. und 21. Januar zusammen mit der vom 11. Februar, die er beibehalten wollte, damit der Zeitraum möglichst groß blieb. Diese Rechnung ergab für ihn bei der vierten Hypothese die folgenden Elemente:
\begin{center}
{\small\setlength{\tabcolsep}{2pt}
\begin{tabular}{l r}
Aphel $\ .\ .\ .\ .\ .\ .\ .\ .\ .$&$330^0\ 33'\ 20''$\\
$\AscendingNode\ .\ .\ .\ .\ .\ .\ .\ .\ .\ .\ .\ .\ .$&$81\quad2\ 35$\\
Bahnneigung $\ .\ .\ .\ .\ .\ .\ .$&$10\ 36\ 30$\\
Logarithmus der großen Halbachse&$0.4370335$\\
Exzentrizität $\ .\ .\ .\ .\ .\ .\ .\ .\ .$&$0.0705553$\\
Epoche wie oben $\ .\ .\ .\ .\ .\ .\ .$&$76^0\ 28'\ 14''.27$\\[4pt]
\multicolumn{2}{c}{Hieraus folgt:}\\[3pt]
größte Mittelpunktsgleichung $\ .\ .\ .$&$8^0\ 5'\ 19''$\\
große Halbachse $\ .\ .\ .\ .\ .\ .\ .\ .$&$2.73548$\\
siderische Umlaufzeit $\ .\ .\ .\ .\ .\ .$&$1652\frac12$ Tage\\
mittlere tägliche siderische Bewegung $\ .$&$784''.25$
\end{tabular}}
\end{center}
Mit diesen Elementen stimmen sämtliche im Septemberheft auf S.\ 280 angeführten Beobachtungen \textsc{Piazzi}s wie folgt überein:
\begin{center}
{\small\setlength{\tabcolsep}{4pt}
\begin{tabular}{r|r|r|r|r}
1801&\multicolumn{2}{c|}{Berechnete}&\multicolumn{2}{c}{Fehler der}\\
&Länge&Breite&Länge&Breite\\\hline
Jan.\ 1&$53^0\ 22'\ 58''42$&$3^0\ 6'\ 42''09$&$+0''12$&$-0''01$\\
2&$53\ 19\ 37.02$&$3\ 2\ 23.78$&$-7.28$&$-1.12$\\
$\mathord{:}\mathord{:}$ 3&$53\ 16\ 43.67$&$2\ 58\ 6.70$&$-14.93$&$-3.20$\\
4&$53\ 14\ 14.03$&$2\ 53\ 51.16$&$-1.47$&$-4.44$\\
$\mathord{:}\mathord{:}$ 10&$53\ 7\ 54.51$&$2\ 28\ 53.13$&$-4.59$&$-7.47$\\
13&$53\ 10\ 18.77$&$2\ 16\ 48.78$&$-18.83$&$-10.92$\\
14&$53\ 11\ 55.25$&$2\ 12\ 51.23$&$-5.95$&$-5.47$\\
19&$53\ 26\ 0.37$&$1\ 53\ 34.19$&$+1.17$&$-4.01$\\
21&$53\ 34\ 22.68$&$1\ 46\ 5.98$&$+1.38$&$-0.02$\\
22&$53\ 39\ 7.88$&$1\ 42\ 25.06$&$+6.08$&$-3.04$
\end{tabular}
\par\medskip
\begin{tabular}{r|r|r|r|r}
1801&\multicolumn{2}{c|}{Berechnete}&\multicolumn{2}{c}{Fehler der}\\
&Länge&Breite&Länge&Breite\\\hline
Jan.\ 23&$53^0\ 44'\ 15''74$&$1^0\ 38'\ 46''25$&$+0''04$&$-5''85$\\
28&$54\ 15\ 20.88$&$1\ 21\ 4.07$&$+5.18$&$-2.83$\\
30&$54\ 30\ 4.25$&$1\ 14\ 13.73$&$+5.25$&$-2.27$\\
31&$54\ 38\ 11.25$&$1\ 10\ 51.99$&$+3.95$&$-2.61$\\
$\mathord{:}\mathord{:}$ Feb.\ 1&$54\ 46\ 28.37$&$1\ 7\ 32.09$&$+9.07$&$+1.19$\\
2&$54\ 55\ 5.09$&$1\ 4\ 14.25$&$+7.19$&$+3.75$\\
5&$55\ 22\ 50.25$&$0\ 54\ 32.88$&$+6.85$&$+3.98$\\
8&$55\ 53\ 23.15$&$0\ 45\ 9.20$&$-6.38$&$+4.20$\\
11&$56\ 26\ 39.97$&$0\ 36\ 2.90$&$-0.03$&$+0.00$
\end{tabular}}
\end{center}
\SourceLeaf{206}%
So ausgezeichnet diese Übereinstimmung ist, hält es Dr.\ \textsc{Gauss} selbst doch für gut möglich, dass seine Elemente von den wahren beträchtlich stärker abweichen als voneinander, da der durchlaufene Teil der Bahn so klein ist und von der ersten bis zur letzten Beobachtung nur $9^0\ 15'\ 35''$ beträgt. Da diese Bahn aber immerhin \textit{möglich}, wenn nicht \textit{höchst wahrscheinlich} ist, haben beobachtende Astronomen doppelten Grund, sie bei der Suche nach diesem Gestirn zu berücksichtigen und sich vor allem an ihr zu orientieren. Von allen bisher erprobten und berechneten Bahnen bringt keine sämtliche Beobachtungen so eng und genau zur Übereinstimmung wie gerade diese \textsc{Gauss}'sche Ellipse.

Wir bemerken außerdem, dass Dr.\ \textsc{Gauss} bei all diesen Berechnungen sowohl den im Oktoberheft auf S.\ 365 angezeigten Druckfehler von einer Minute im Ort der Sonne als auch einen zweiten am 13. Jan. berücksichtigt hat, der ebenfalls die Länge der Sonne betraf. Für diese letzte Länge hat er $9^{\mathrm z}\ 23^0\ 13'\ 13.''8$ angenommen.

So weit war Dr.\ \textsc{Gauss} mit seinen Berechnungen der Bahn der \textit{Ceres Ferdinandea} gekommen, als er sie uns freundlicherweise mitteilte. Inzwischen hatten wir die Abhandlung \textsc{Piazzi}s samt seinen berichtigten Beobachtungen erhalten und im Novemberheft einen getreuen Auszug daraus mitgeteilt. Diese Änderungen waren gering, am beträchtlichsten war die um $15''$ verminderte Rektaszension vom 11. Febr., und sie änderten die vorigen Elemente nur wenig. Doch hatten sich hier und da einige bedeutendere Druck- und Rechenfehler in die Reduktion dieser Beobachtungen eingeschlichen. Da sich die Fehler leicht beseitigen und die Elemente daraufhin bald verbessern ließen, teilten wir Dr.\ \textsc{Gauss} zu diesem Zweck sämtliche Beobachtungen \textsc{Piazzi}s in der italienischen Originalschrift mit. Noch bevor er unsere Zuschrift erhielt, hatte er jedoch bereits mit der Berechnung einer neuen Bahn begonnen. So gering die Abweichungen seiner zweiten Elemente von den Beobachtungen schon waren, wiesen sie doch noch eine auffallende Regelmäßigkeit auf. Daraus ließ sich mit Gewissheit voraussehen, dass sich die Fehler noch beträchtlich weiter verringern ließen. Er machte den Versuch und war mit der Arbeit bereits weitgehend fertig, als er die ihm zugesandten berichtigten Beobachtungen erhielt, so wie \textsc{Piazzi} selbst sie reduziert und berechnet hatte. Daher hielt er es für richtig, die Arbeit zunächst ganz abzuschließen, zumal er im Voraus wusste, dass die Länge für den 11. Febr. nach den neuen Elementen um $6''$ kleiner ausfallen und sich damit von selbst schon fast um die Hälfte dem berichtigten Beobachtungswert nähern würde, der ihm bis dahin natürlich unbekannt war. Auf diese Weise fand er die folgenden \textit{dritten} Elemente:
\begin{center}
{\small\setlength{\tabcolsep}{2pt}
\begin{tabular}{l r}
Aphel $\ .\ .\ .\ .\ .\ .\ .\ .\ .\ .$&$326^0\ 53'\ 50''$\\
$\AscendingNode\ .\ .\ .\ .\ .\ .\ .\ .\ .\ .\ .\ .\ .$&$81\quad1\ 44$\\
Bahnneigung $\ .\ .\ .\ .\ .\ .\ .\ .$&$10\ 36\ 21$\\
Logarithmus der großen Halbachse $\ .$&$0.4414902$\\
Exzentrizität $\ .\ .\ .\ .\ .\ .\ .\ .\ .\ .$&$0.0819603$\\
Epoche $\ .\ .\ .\ .\ .\ .\ .\ .\ .\ .\ .\ .$&$77^0\ 34'\ 28''$\\[4pt]
\multicolumn{2}{c}{Hieraus ferner:}\\[3pt]
Größte Mittelpunktsgleichung $\ .\ .\ .$&$9^0\ 23'\ 57''$\\
Mittlere tägliche siderische heliozentrische Bewegung&$772''275$\\
Tägliche tropische heliozentrische Bewegung $\ .$&$772.413$\\
Tropische Umlaufzeit $\ .\ .\ .\ .\ .\ .$&$1677.8$ Tage
\end{tabular}}
\end{center}
Diese Elemente stellen die von uns reduzierten und im Septemberheft auf S.\ 280 angeführten Beobachtungen \textsc{Piazzi}s folgendermaßen dar:
\expandafter\def\csname HistoricalNoteText17\endcsname{In den Originalbeobachtungen \textsc{Piazzi}s auf S.\ 565 des Novemberhefts sind die Längen der Sonne am 13. Januar und 8. Februar um eine Minute zu groß. Alle geozentrischen Breiten sind um etwa $10''$ zu klein. Das kommt zweifellos daher, dass \textsc{Piazzi} bei der Reduktion seiner Beobachtungen eine andere Schiefe der Ekliptik verwendet hat, wahrscheinlich die von $23^0\ 28'\ 1''$ aus der \textit{Conn.\ d.\ tems Année} IX.}
\SourceLeaf{207}%
\begin{center}
{\small\setlength{\tabcolsep}{4pt}
\begin{tabular}{r|r|r|r|r}
1801&\multicolumn{2}{c|}{Berechnete}&\multicolumn{2}{c}{Fehler der}\\
&Länge&Breite&Länge&Breite\\\hline
Jan.\ 1&$53^0\ 23'\ 2''34$&$3^0\ 6'\ 43''63$&$+4''04$&$+1''53$\\
2&$53\ 19\ 41.24$&$3\ 2\ 25.68$&$-3.06$&$+0.78$\\
3&$53\ 16\ 48.05$&$2\ 58\ 8.97$&$-10.35$&$-0.93$\\
4&$53\ 14\ 18.47$&$2\ 53\ 53.79$&$+2.97$&$-1.81$\\
10&$53\ 7\ 58.37$&$2\ 28\ 57.12$&$-0.73$&$-3.48$\\
13&$53\ 10\ 21.60$&$2\ 16\ 52.89$&$-16.00$&$-6.81$\\
14&$53\ 11\ 57.70$&$2\ 12\ 55.36$&$-3.50$&$-1.34$\\
19&$53\ 26\ 0.59$&$1\ 53\ 38.01$&$+1.39$&$-0.19$\\
21&$53\ 34\ 21.99$&$1\ 46\ 9.53$&$+0.69$&$+3.53$\\
22&$53\ 39\ 6.69$&$1\ 42\ 28.45$&$+4.89$&$+0.35$
\end{tabular}
\par\medskip
\begin{tabular}{r|r|r|r|r}
1801&\multicolumn{2}{c|}{Berechnete}&\multicolumn{2}{c}{Fehler der}\\
&Länge&Breite&Länge&Breite\\\hline
Jan.\ 23&$53^0\ 44'\ 14''08$&$1^0\ 38'\ 49''44$&$-1''62$&$-2''66$\\
28&$54\ 15\ 17.11$&$1\ 21\ 5.91$&$+1.41$&$-0.99$\\
30&$54\ 30\ 9.76$&$1\ 14\ 15.12$&$+0.76$&$-0.88$\\
31&$54\ 38\ 6.44$&$1\ 10\ 52.81$&$-0.86$&$-1.79$\\
Febr.\ 1&$54\ 46\ 23.22$&$1\ 7\ 32.54$&$+3.92$&$+1.64$\\
2&$54\ 54\ 59.71$&$1\ 4\ 14.30$&$+1.81$&$+3.80$\\
5&$55\ 22\ 44.30$&$0\ 54\ 31.72$&$+0.90$&$+2.82$\\
8&$55\ 53\ 17.01$&$0\ 45\ 6.65$&$-12.49$&$+1.63$\\
11&$56\ 26\ 34.10$&$0\ 35\ 58.96$&$-5.90$&$-3.94$
\end{tabular}}
\end{center}
Die von \textsc{Piazzi} angezeigten Fehler und Berichtigungen ergeben nun folgende Änderungen. Wegen der um $15''$ verminderten Rektaszension vom 11. Febr. erhält man bei einer Schiefe der Ekliptik von $23^0\ 28'\ 12''$ die Länge $56^0\ 26'\ 26''1$ und die Breite $35'\ 59''7$. Damit beträgt der Längenfehler $+8''0$ und der Breitenfehler $-0''74$. Beim Nachrechnen der Reduktion von Rektaszension und Deklination für diejenigen Beobachtungen, bei denen \textsc{Piazzi}s Angaben beträchtlich von unseren abweichen, erhielt Dr.\ \textsc{Gauss} für den 8. Febr. die Länge $55^0\ 53'\ 17''7$. Der Fehler beträgt folglich $-0''7$\HistoricalNote{17}. Im Übrigen kommt Dr.\ \textsc{Gauss} zu denselben Ergebnissen wie wir, selbst bei der Länge am 13. Januar, die bei \textsc{Piazzi} so stark abweicht. Gerade diese Beobachtung weicht sowohl von den letzten als auch von den vorigen Elementen am stärksten ab und lässt sich nur auf Kosten der übrigen Beobachtungen genau darstellen. Daher vermutet Dr.\ \textsc{Gauss}, dass dabei irgendwo ein kleines Versehen unterlaufen ist. Im Übrigen sieht man leicht, dass die übrigen von \textsc{Piazzi} angegebenen kleinen Berichtigungen (Novemberheft, S.\ 573) die Genauigkeit, mit der die Beobachtungen durch diese letzten Elemente dargestellt werden, wenig verändern.

So gering jedoch alle Fehler sind, hat Dr.\ \textsc{Gauss} noch einen vierten Versuch unternommen, die Genauigkeit etwas weiter zu steigern. Bei diesen letzten Elementen sind die Längenfehler im Februar sämtlich positiv, wenn man auch die kleinen Berichtigungen \textsc{Piazzi}s berücksichtigt. Auch die Breitenfehler liegen gegen Ende überwiegend auf einer Seite. Dr.\ \textsc{Gauss} hat daher neue Elemente gesucht, nach denen die berechneten Längen und Breiten im Februar etwas kleiner werden, während die übrigen ungefähr gleich bleiben. Obwohl er diese Elemente nicht mit sämtlichen Beobachtungen verglichen hat und dies zum Teil auch ganz überflüssig ist, glaubt er doch im Voraus versichern zu können, dass sie zwischen den Beobachtungen so weit wie möglich die Mitte halten und dass die Fehler bei jeder Beobachtung (außer denen vom 3. und 13. Januar) in Länge und Breite innerhalb von $5''$ bleiben, keinerlei Regelmäßigkeit mehr zeigen und ihre Vorzeichen ganz unregelmäßig wechseln. Dr.\ \textsc{Gauss} bezweifelt mit gutem Grund, dass man \textit{bei genau denselben Daten} mit anderen Elementen eine \textit{merklich größere} Übereinstimmung erzielen könnte. Damit will er jedoch durchaus nicht sagen oder so verstanden werden, dass andere, erheblich verschiedene Elemente keine ähnliche Übereinstimmung ergeben könnten, zumal wenn man bei dieser empfindlichen Rechnung, bei der schon wenige Sekunden einen merklichen Unterschied machen, etwas andere Bestimmungen für die Sonne verwendete. So unterscheidet sich die Länge der Sonne bei \textsc{Piazzi} im Februar um eine halbe Minute von unserer, obwohl gerade hier in der Nähe der Quadratur der Einfluss der Sonnenlänge geringer ist als in anderen Lagen. Dr.\ \textsc{Gauss} hält es daher für möglicherweise nützlich, die Fehler der
\SourceLeaf{208}%
Sonnentafeln für diese Zeiten anhand sehr genauer Beobachtungen zu bestimmen und die Orte der Sonne danach zu verbessern. Diese \textit{vierten} Elemente lauten nun wie folgt:
\begin{center}
{\small\setlength{\tabcolsep}{2pt}
\begin{tabular}{l r}
Aphel $\ .\ .\ .\ .\ .\ .\ .\ .\ .\ .$&$326^0\ 27'\ 38''$\\
$\AscendingNode\ .\ .\ .\ .\ .\ .\ .\ .\ .\ .\ .\ .\ .$&$81\quad0\ 44$\\
Neigung $\ .\ .\ .\ .\ .\ .\ .\ .\ .\ .\ .$&$10\ 36\ 57$\\
Logarithmus der großen Halbachse $\ .$&$0.4420527$\\
Exzentrizität $\ .\ .\ .\ .\ .\ .\ .\ .\ .\ .$&$0.0825017$\\
Epoche 1800 31. Dez.\ $\ .\ .\ .\ .\ .\ .\ .$&$77^0\ 36'\ 34''$\\[4pt]
\multicolumn{2}{c}{Hieraus:}\\[3pt]
größte Mittelpunktsgleichung $\ .\ .\ .\ .$&$9^0\ 27'\ 41''$\\
\multicolumn{2}{l}{tägliche mittlere heliozentrische tropische Bewegung}\\
$\ .\ .\ .\ .\ .\ .\ .\ .\ .\ .\ .\ .$&$770''914$
\end{tabular}}
\end{center}
Aus diesen Elementen hat Dr.\ \textsc{Gauss} die folgenden Orte der \textit{Ceres Ferdinandea} vorausberechnet. Die Zeit ist mittlere Zeit für Mitternacht in \textit{Palermo}.
\begin{center}
{\fontsize{9}{11}\selectfont\setlength{\tabcolsep}{3pt}
\begin{tabular}{r|r|r|r|r|r}
1801&\shortstack{Geozentrische\\Länge}&\shortstack{Geozentrische\\Breite\\nördlich}&\shortstack{Logarithmus\\des Abstands\\von der $\EarthSymbol$}&\shortstack{Logarithmus\\des Abstands\\von der $\odot$}&\shortstack{Verhältnis\\der sichtbaren\\Helligkeit}\\\hline
November 25&$5^{\mathrm z}\ 20^0\ 16'$&$9^0\ 25'$&$0.42181$&$0.40468$&$0.6102$\\
Dezember 1&$5\ 22\ 15$&$9\ 48$&$0.40940$&$0.40472$&$0.6459$\\
7&$5\ 24\ 7$&$10\ 12$&$0.39643$&$0.40479$&$0.6855$\\
13&$5\ 25\ 51$&$10\ 37$&$0.38296$&$0.40488$&$0.7290$\\
19&$5\ 27\ 27$&$11\ 4$&$0.36902$&$0.40499$&$0.7770$\\
25&$5\ 28\ 53$&$11\ 32$&$0.35468$&$0.40512$&$0.8295$\\
31&$6\ 0\ 10$&$12\ 1$&$0.34000$&$0.40528$&$0.8869$
\end{tabular}}
\end{center}
\Needspace{12\baselineskip}
Falls man den Ort des Planeten nach diesen Elementen genauer oder für eine längere Zeit berechnen möchte, geben wir dafür hier noch folgende Formeln:

\begin{center}1) Zur Berechnung der Mittelpunktsgleichung $=M$\end{center}
\[
\begin{aligned}
M={}&-34005''494\ \mathrm{Sin.\ Anom.\ med.}\\
&+1750''951\ \mathrm{Sin.}\ 2\ \mathrm{A.\ m.}-124''951\ \mathrm{Sin.}\ 3\ \mathrm{A.\ m}\\
&+10''192\ \mathrm{Sin.}\ 4\ \mathrm{A.\ m.}-0''901\ \mathrm{Sin.}\ 5\ \mathrm{A.\ m.}\\
&+0''083\ \mathrm{Sin.}\ 6\ \mathrm{A.\ m.}
\end{aligned}
\]
\begin{center}2) Für den Radiusvektor $=r$\end{center}
\[
\begin{aligned}
r={}&2.776695+0.2373586\ \mathrm{Cosin.\ Anom.\ med.}\\
&-0.0093752\ \mathrm{Cosin.}\ 2\ \mathrm{A.\ m.}\\
&+0.0005789\ \mathrm{Cosin.}\ 3\ \mathrm{A.\ m.}\\
&-0.0000423\ \mathrm{Cosin.}\ 4\ \mathrm{A.\ m.}\\
&+0.0000034\ \mathrm{Cosin.}\ 5\ \mathrm{A.\ m.}\\
&-0.0000003\ \mathrm{Cosin.}\ 6\ \mathrm{A.\ m.}
\end{aligned}
\]
\[
\text{Oder:}\quad r=\frac{7.60570}{2.767278\pm0.2283053\ \mathrm{Cosin.\ Anom.\ ver.}}
\]
\begin{center}3) Für die heliozentrische Breite $=\lambda$.\end{center}
\[
\mathrm{Log.\ Sin.}\ \lambda=9.2653438+\mathrm{Log.\ Sin.\ Arg.\ Latit.}
\]
\begin{center}4) Für die Reduktion der heliozentrischen Länge auf die Erdbahn $=\varepsilon$.\end{center}
\[
\begin{aligned}
\alpha)&\quad\mathrm{Log.\ Tang.}\ \varphi=9.9925025+\mathrm{Log.\ Tang.\ Arg.\ Latit.}\\
\beta)&\quad\varepsilon=\mathrm{Arg.\ Latit.}-\varphi.
\end{aligned}
\]
\begin{center}5) Für die Reduktion des Radiusvektors oder die verkürzte Distanz $=\rho$.\end{center}
\[
\rho=\mathrm{Cosin.\ Latit.\ helioc.}\times\mathrm{Distant.\ vera.}
\]
\begin{center}Oder für den Logarithmus der Verkürzung selbst $=\mathrm{Log.}\ \mu$.\end{center}
\[
\mathrm{Log.}\ \mu=10.000000-\mathrm{Log.\ Cosin.\ Latit.\ helioc.}
\]
\begin{center}6) Für die Aberration des Lichts in Länge, Breite, Rektaszension und Deklination $=a$.\end{center}
\[
\begin{aligned}
\mathrm{Log.}\ a={}&\pm\mathrm{Log.\ Dist.}\ \text{à}\ \EarthSymbol+\mathrm{Log.\ mot.\ horar.\ geoc.}\\
&+\mathrm{Log.}\ 7.751007.
\end{aligned}
\]

\SourceLeaf{209}%
Wir schließen diese Nachricht mit einer Bemerkung von Dr.\ \textsc{Gauss} zur Neigung der Bahn der \textit{Ceres}, deren Größe vielen Astronomen so aufgefallen ist. Er stimmt dem Urteil einiger Astronomen völlig zu, dass uns kein physikalischer Grund dazu berechtigt, bei allen Himmelskörpern unseres Sonnensystems mit annähernd kreisförmiger Bahn auch eine geringe Neigung gegen die Ekliptik zu erwarten. Dennoch scheint ihm, dass sich das Paradoxe daran beträchtlich abschwächen und die Analogie einigermaßen retten lässt, wenn man diese Planetenbahn auf die Ebene bezieht, die ihr als Bezugsebene zukommt, nämlich auf den Sonnenäquator, wie \textsc{La Place} es bei den Bahnen der Uranusmonde getan hat (\textit{A.\ G.\ E.}\ Bd.\ II, S.\ 259). Vergleichen wir die Bahnebenen sämtlicher Planeten miteinander, zeigt sich dann, dass die Bahn der \textit{Ceres} gegen die Bahn keines anderen Planeten so stark geneigt ist wie gegen die Erdbahn. Wenn wir also die Bahn der \textit{Ceres} mit der Erdbahn vergleichen, vergleichen wir gerade die beiden Extreme unseres Sonnensystems. Denkt man sich aber eine Ebene, die ungefähr in der Mitte zwischen den Bahnebenen sämtlicher \textit{acht} Planeten liegt, so sind die Neigungen gegen diese Ebene alle klein genug. Es scheint sehr bemerkenswert, dass gerade der Sonnenäquator ungefähr so liegt, mit der einzigen Einschränkung, dass die Bahn der \textit{Ceres} ebenso wie die der Venus und des Merkur am wenigsten gegen ihn geneigt ist, etwa 3 bis 4 Grad, während die Erdbahn gerade am stärksten gegen ihn geneigt ist.
\begin{flushright}[\textsc{von Zach}.]\end{flushright}


